\chapter{Análisis y resultados}

En este capítulo se presentan los resultados obtenidos durante el diagnóstico realizado en la fábrica Confort Muebles antes de la implementación de la aplicación web. Para ello se emplearon técnicas de investigación cualitativa y cuantitativa descritas en el capítulo anterior, entre ellas la entrevista semiestructurada aplicada al administrador de la empresa y la observación directa de los procesos administrativos.
Posteriormente se presentan los resultados obtenidos durante la implementación del sistema, con la finalidad de analizar el impacto de la aplicación web en la automatización del proceso de gestión de proformas, clientes, inventario y pagos.
\section{Resultados del diagnóstico inicial}
Como parte del diagnóstico inicial se aplicó una entrevista semiestructurada al administrador de Confort Muebles, quien constituye el único usuario directo del sistema y el responsable de la gestión administrativa de la empresa. La entrevista permitió identificar las principales dificultades presentes en los procesos actuales, así como los requerimientos funcionales considerados para el desarrollo de la aplicación web.
\subsection{Resultados de la entrevista al administrador}
\begin{table}[H]
\centering
\caption{Resultados de la entrevista al administrador de Confort Muebles}
\label{tab:entrevista}

\small
\renewcommand{\arraystretch}{1.3}

\begin{tabular}{p{6.5cm} p{8cm}}
\toprule
\textbf{Pregunta} & \textbf{Respuesta resumida} \\
\midrule

¿Cuánto tiempo tarda en elaborar una proforma? &
Entre 20 y 40 minutos, dependiendo de la complejidad del pedido. \\

¿Ha presentado errores durante la elaboración de las cotizaciones? &
Sí. Principalmente en precios, materiales y cantidades. \\

¿Qué información consulta con mayor frecuencia? &
Clientes, cotizaciones, inventario y productos. \\

¿Considera necesario implementar un sistema web? &
Sí, porque permitirá reducir errores, agilizar los procesos y mejorar el control de la información. \\

¿Qué mejoras espera obtener con el sistema? &
Reducir el tiempo de elaboración de proformas, disminuir errores y mejorar la organización administrativa. \\

\bottomrule
\end{tabular}

\vspace{0.3cm}
\footnotesize
\textit{Nota.} Elaboración propia a partir de la entrevista realizada al administrador de Confort Muebles.

\end{table}
\subsubsection*{Análisis de los resultados}

Los resultados obtenidos mediante la entrevista evidencian que los procesos administrativos de Confort Muebles continúan desarrollándose principalmente de forma manual. El administrador indicó que los pedidos son recibidos de manera presencial, mediante llamadas telefónicas y por WhatsApp, por lo que la información debe registrarse posteriormente en documentos físicos o archivos independientes.

Asimismo, se determinó que la elaboración de una proforma requiere entre veinte y cuarenta minutos, debido a que el cálculo del precio de venta se realiza manualmente considerando materiales, mano de obra y porcentaje de utilidad. Esta metodología incrementa el tiempo de atención al cliente y aumenta la probabilidad de errores durante el proceso de cotización.

Durante la entrevista también se identificó que la empresa ha presentado pérdidas ocasionales de información y errores relacionados con precios, cantidades y materiales incluidos en las cotizaciones. Estas dificultades son consecuencia de la inexistencia de una base de datos centralizada que permita administrar de manera organizada la información de clientes, productos, inventario y ventas.

El administrador manifestó además la necesidad de disponer de reportes administrativos sobre ventas, cotizaciones, inventario, clientes y pagos pendientes, así como la posibilidad de acceder a la información desde cualquier dispositivo con conexión a Internet. Estas funcionalidades fueron consideradas durante el desarrollo de la aplicación web.
\subsubsection*{Interpretación de los resultados}

La información obtenida permitió identificar las principales necesidades administrativas de Confort Muebles y justificar el desarrollo de una aplicación web orientada a la automatización del proceso de gestión de proformas.

Los resultados evidencian que la dependencia de registros manuales genera retrasos, dificulta el seguimiento de clientes y aumenta la probabilidad de errores administrativos. Asimismo, se determinó que la centralización de la información y la automatización de los procesos constituyen necesidades prioritarias para mejorar la eficiencia operativa de la empresa.

En función de estos resultados se diseñó una aplicación web que incorpora módulos para la administración de clientes, productos, inventario, cotizaciones, pedidos y pagos, además de un panel principal con indicadores estadísticos y reportes administrativos que facilitan la toma de decisiones y optimizan el control de la información.
